\documentclass[literature.tex]{subfiles}

\begin{document}
\prose{Smrť krásných srnců}{Ota Pavel}
\teorieA{
epická próza, autobiografický soubor povídek (vzpomínková próza s~prvky tragikomiky)
}{}{
Jednoduchý, hovorový a~lidový jazyk s~prvky rybářského slangu, přímou řečí a~humoristickými obraty. Časté popisy přírody (Berounka, Křivoklátsko) mají poetické vyznění. Kontrast lehkého, humorného tónu předválečných příběhů s~vážnými válečnými motivy.
}{
Přímá řeč, opakování, ironie, humoristické zkratky a~tragikomické pointy.
}{
Přirovnání (\texthl{štiky jak krokodýlové}), metafora (srnci jako symbol krásy a~nevinnosti ničené válkou), personifikace přírody a~ryb (kapři jako „děti“ tatínka).
}{
Vyprávění v~\wemph{ich-formě} z~pohledu malého Oty (později dospívajícího chlapce). Chronologický postup s~retrospektivami. Dominuje vyprávěcí a~popisný styl spojený s~humoristickým a~nostalgickým tónem. Tragikomický kontrast idylického dětství a~válečných útrap.
}{
Tenkrát jsem nevěděl proč, ale dnes to vím. Můj tatínek pochopil už tenkrát, že jednou můžu uvidět bulváry Paříže a~mrakodrapy New Yorku, ale už nikdy nebudu moci týdny bydlet v~chalupě, kde voní v~peci chleba a~v~máselnici se tluče máslo...\\
Kdysi tatínek hledal ten kraj po~čuchu. Ve třicátých letech minul s~naším osobním šoférem Tondou Valentou hrad Křivoklát a~stoupal strašnými serpentinami dále na~západ podél řeky Berounky, ve které tenkrát plavaly štiky jak krokodýlové...
}
\teorieB{
\wtem{Tatínek}{ústřední postava, charismatický, podnikavý židovský obchodník (Electrolux), snílek, vášnivý rybář a~dobrodruh plný vitality a~humoru. Nezdolný optimista (\texthl{„Vyděláme na tom majlant!“}), ochotný riskovat život pro rodinu. Symbol lidské statečnosti a~lásky k~přírodě.}
\wtem{Vypravěč (Ota Pavel)}{malý chlapec (autor v~ich-formě), citlivý pozorovatel, obdivuje otce a~lásku k~přírodě.}
\wtem{Maminka}{tolerantní, laskavá křesťanka, opora rodiny, odpouští tatínkovy avantýry a~výstřelky.}
\wtem{Strejda Prošek}{pytlák z~Luhu pod Bánovem (u~Křivoklátu), „král pytláků“, všeuměl sžitý s~přírodou.}
\wtem{Holan}{divoký vlčák, vycvičený na~pytlačení, věrný společník Proška.}
\wtem{Bratři Jiří a~Hugo}{starší bratři, odvlečeni do~koncentračního tábora.}
}{
Soubor 7–9 povídek chronologicky vzpomíná na~autorovo dětství a~rodinný život smíšené židovské rodiny. Předválečné idylické zážitky (rybaření, podnikání, příroda) kontrastují s~protektorátní perzekucí a~válečnými útrapami. Hlavní motivy: láska k~přírodě (Berounka, ryby, srnci), rodinná soudržnost, otcova obětavost a~nezdolnost, tragikomika života.\\
\vspace{1pt}\\
\textbl{Nejdražší ve střední Evropě}: Tatínek koupí „rybník plný ryb“, ale najde v~něm jen jednoho kapra. Pomstí se bývalému majiteli prodejem vadné ledničky.\\
\textbl{Ve službách Švédska}: Tatínkova podnikavost u~Electroluxu – stane se mistrem prodeje vysavačů a~ledniček díky šarmu a~tajnému obdivu k~paní Irmě.\\
\textbl{Smrť krásných srnců} (nejznámější, zfilmována): Za protektorátu tatínek riskuje život (Davidova hvězda, zákaz pohybu) a~s~Proškovým psem Holanem uloví srnce, aby nasytil syny před transportem do~koncentráku. Srnci symbolizují krásu a~nevinnost ničenou válkou.\\
\textbl{Kapři pro Wehrmacht} (zfilmována): Tatínkovi zaberou rybník Němci. On tajně vyloví své „kapříky“ a~Němci najdou prázdný rybník – malý akt odporu.\\
\textbl{Jak jsme se střetli s~Vlky}: Krátká humorná příhoda o~střetu s~„Vlky“ (pravděpodobně sousedy nebo úřady).\\
\textbl{Otázka hmyzu vyřešena}: Poválečné podnikání s~mucholapkami BOMBA-CHEMIK – rychlý úspěch i~pád (mouchy „nesmrtelné“).\\
\textbl{Králíci s~moudrýma očima}: Poslední štace rodiny u~Radotína – chov šampaňských králíků, idylický, ale tragický konec (tatínkova smrt).
}{
Soubor samostatných, ale tematicky propojených povídek v~chronologickém pořadí (předválečné idylické období → protektorát a~válka → poválečné roky). Ich-forma. Prostor převážně Buštěhrad, Křivoklátsko, řeka Berounka a~okolí. Čas: 30.–50./60. léta 20. století.
}{
Nostalgická vzpomínka na~dětství, lásku k~přírodě a~rodině. Oslava otcovy statečnosti, optimismu a~lidskosti v~těžkých časech. Kritika války a~rasové perzekuce Židů osobním, nepatetickým způsobem. Motivy: vztah člověka k~přírodě, rodinná soudržnost, tragikomika života, obětavost rodičů a~ztráta idyly (\texthl{smrt krásných srnců} jako metafora války).
}\newpage
\historie{
Dílo vzniklo v~době normalizace (1971). Odráží 2. světovou válku, protektorát, perzekuci Židů, transporty do~koncentračních táborů a~poválečné období. Kontrast předválečné svobody s~válečným útlakem.
}{}{}{}{}{
Vlastním jménem Otto Popper (1930 Praha – 1973 Praha). Český prozaik, novinář a~sportovní reportér židovského původu. Za~války přežil s~matkou (otec a~bratři v~táborech). Pracoval jako sportovní redaktor. V~roce 1964 propukla duševní choroba (maniodepresivní psychóza). Dílo psal často v~léčebně jako terapii vzpomínkami na~„nejšťastnější okamžiky“. Zemřel ve~42 letech.
}{
Mezi nejznámější díla patří \textsc{Jak jsem potkal ryby} (navazuje tematicky), \textsc{Zlatí úhoři}, sportovní prózy jako \textsc{Dukla mezi mrakodrapy} nebo \textsc{Pohádka o~Raškovi}.
}
\kritika{}{}{
Témata lásky k~přírodě, rodinné soudržnosti, odvahy v~těžkých časech a~osobní vzpomínky na~holokaust zůstávají nadčasová. Dílo připomíná válečné útrapy nepateticky a~lidsky.
}
\nazor{
Ota Pavel proměnil vzpomínky na~dětství a~charismatického tatínka v~dojemný, humorný i~dojemný soubor povídek plný lásky k~přírodě a~rodině. Nejsilnější je titulní povídka o~lovu srnce – tatínek riskuje život, aby synům dal sílu před transportem. Tragikomický kontrast (smích před válkou vs. hrůza okupace) funguje výborně. Dílo je čtivé, nostalgické a~lidské. Rozhodně patří k~nejdojemnějším českým vzpomínkovým prózám – připomíná, jak krásný může být obyčejný život i~v~nejtěžších časech.
}
%\explain{}
\wpage
\end{document}